\section{The final handoffs and the remaining construction theorem}
\label{ex-assembly}

The final assembly must use one hierarchy throughout. An operator-product
expansion for one subsequential limit and a reconstruction for another would
not establish the desired result. We first state the remaining two properties
of that common hierarchy, then isolate precisely what is needed to assemble it.

\subsection{The operator-product and trace requirements}
Let \(\boldsymbol A=(A_i)_{i\in I}\) and
\(\boldsymbol B=(B_j)_{j\in J}\) be finite lists of graded fields and
write \(w(\boldsymbol A)=\sum_{i\in I}\deg A_i\). For each grade
\(\ell\), let \((E_{\ell,\alpha})_{\alpha\in\mathcal C_\ell}\)
be a finite collection of fields. The coefficients
\(c_{\boldsymbol A}^{\ell,\alpha}(\boldsymbol x;o)\) do not depend
on the phase. The cluster geometry at scale \(r>0\) is
\[
 |x_i-o|\le Mr,\qquad |x_i-x_{i'}|\ge\eta r\quad(i\ne i'),
\]
where \(M\ge0\) and \(\eta>0\). Spectators satisfy
\(|y_j-o|\ge c\) and \(|y_j-y_{j'}|\ge c\) for \(j\ne j'\),
where \(c>0\). For every \(\epsilon>0\), each fixed coefficient obeys
\begin{equation}\label{ex-ope-coefficient}
 |c_{\boldsymbol A}^{\ell,\alpha}(\boldsymbol x;o)|
 \le r^{\ell-w(\boldsymbol A)-\epsilon}
\end{equation}
at sufficiently small \(r\), uniformly in the indicated cluster geometry.
For each fixed truncation degree \(D\), the expansion requires
\begin{align}
 \bigg|S^s_{\boldsymbol A,\boldsymbol B}(\boldsymbol x,\boldsymbol y)
 -\sum_{\ell=0}^D\sum_{\alpha\in\mathcal C_\ell}
 c_{\boldsymbol A}^{\ell,\alpha}(\boldsymbol x;o)
 S^s_{E_{\ell,\alpha},\boldsymbol B}(o,\boldsymbol y)\bigg|
 \le r^{D+1-w(\boldsymbol A)-\epsilon}.
 \label{ex-ope-remainder}
\end{align}
Here \(s\) can be any character or the mixed state. The small-radius
threshold for the remainder is uniform over these choices and over the
allowed positions, but may depend on the fields, \(D,M,\eta,c,\epsilon\).
There is no assertion of one threshold working for all truncation degrees.

The trace requirement concerns a centered grade-four field \(\Theta\).
Set
\[
 T(r)=S^{\chi_*}_{\Theta,\Theta}
            ((\boldsymbol0,-r/2),(\boldsymbol0,r/2)).
\]
Its vacuum one-point functional vanishes, and the required asymptotic is
\begin{equation}\label{ex-trace-limit}
 \lim_{r\downarrow0}r^8\log^2(r^{-1})T(r)
       =\frac{3\dim_{\mathbb R}\mathfrak g}{\pi^4}>0.
\end{equation}
Theorem~\ref{thm:profile-exclusion} shows that this asymptotic implies the
required exclusion of finite complex linear combinations of mass-shell
tracks, including empty products. The class of profiles and their exact
normalization are given in Section~\ref{sec:trace-profiles}.
Theorem~\ref{thm:trace-asymptotic} supplies the limiting argument from
the trace comparison estimates; producing those estimates remains necessary.

\subsection{A precise conditional assembly}
An internal extraction order is a pair \((u,p)\) of strictly increasing
maps \(\mathbb N\to\mathbb N\). Refining by \((v,q)\) gives
\((u\circ v,p\circ q)\). The full mesh order is \(u(N)=N\).
The need for a full mesh limit at the public endpoint is stronger than
convergence just along \(u\circ v\).

\begin{lemma}[Identifying the full mesh limit]\label{ex-full-mesh}
Fix a test and its insertion labels. Suppose that at every suitable period
\(p\) the full sequence \(W_{N,p}\) converges to some \(\ell_p\).
Suppose also that an ordered extraction supplies numbers \(w_p\) such that
\(W_{u(j),p}\to w_p\), where \(u\) is strictly increasing, and
\(w_{p(q(k))}\to s\). Then \(\ell_p=w_p\) at every suitable period,
and the same volume limit \(s\) holds with the full mesh limits.
\end{lemma}
\begin{proof}
Strict increase implies \(u(j)\ge j\) by induction, so \(u(j)\to\infty\).
The full mesh convergence therefore implies
\(W_{u(j),p}\to\ell_p\). To see uniqueness explicitly, for any
\(\epsilon>0\), take \(j\) sufficiently large that this term is within
\(\epsilon/2\) of each proposed limit. The triangle inequality gives
\(|\ell_p-w_p|<\epsilon\), hence equality. Suitability holds eventually
on the volume sequence, so the two sequences of row limits eventually agree
and have the same limit.
\end{proof}

This lemma explains how a full ultraviolet convergence result can be combined
with a hierarchy constructed by subsequence extraction. It does not obtain
full ultraviolet convergence from subsequential convergence: that is its
additional hypothesis.

\begin{assumption}[Joint construction input]\label{ex-joint-input}
For the fixed group \(G\) and each chosen positive blocking penalty
\(b\), there is one lattice family with the following properties.
Its insertions are continuous, gauge-invariant and uniformly local, and its
bare coefficients satisfy \eqref{ex-running-input}.
For every extraction order \((u,p)\), there is a refinement \((v,q)\)
and one set of graded fields, phase data, correlations and selected
relativistic theory as in Definition~\ref{ex-reconstructed}, for which
all of the following hold simultaneously:
\begin{enumerate}
\item The mixed correlations satisfy \eqref{ex-row-target} along the
\emph{full} mesh sequence at every suitable period and
\eqref{ex-volume-target} along \(p\circ q\).
\item The phase correlations have the stated smoothness, symmetry,
continuity and mixture identity \eqref{ex-mixture}.
\item The fields and coefficients satisfy
\eqref{ex-ope-coefficient}--\eqref{ex-ope-remainder} and the centered
field satisfies \eqref{ex-trace-limit}.
\end{enumerate}
Every assertion refers to the same lattice family and the same selected
hierarchy on that refinement.
\end{assumption}

\begin{proposition}[Assembly from the joint input]\label{ex-conditional-existence}
Input~\ref{ex-joint-input} implies the full existence target stated in
Section~\ref{ex-statement}, including the operator-product and trace
properties above.
\end{proposition}
\begin{proof}
Fix \(G\) and choose \(b=1\). Select the family supplied by
Input~\ref{ex-joint-input} once. Its regularity is part of that input and
is unaffected by any subsequent choice of volume sequence. Now let
\(p:\mathbb N\to\mathbb N\) be an arbitrary strictly increasing
sequence. Apply the input to the extraction order \((\mathrm{id},p)\).
It yields \((v,q)\) and one hierarchy with all three simultaneous
properties. Take \(q\) as the public volume subsequence and retain this
hierarchy. Its ordered Wilson convergence uses the full mesh sequence by
item~1, so no assertion about a mesh subsequence has been substituted for
\eqref{ex-row-target}. Item~2 gives the phase and mixture conditions;
the selected theory is part of the same data. Item~3 supplies the remaining
two properties, with finite-profile exclusion following from
Theorem~\ref{thm:profile-exclusion}. The construction of the family preceded the choice of
\(p\), while \(q\) and the hierarchy were chosen after \(p\), exactly
as required.
\end{proof}

The proposition is only the final logical assembly. The main outstanding
work is to establish Input~\ref{ex-joint-input} from the Wilson measure.
In particular, the finite-dimensional normalization proof does not supply
the estimates uniform in spacing and volume; a convergent mesh subsequence
does not establish the full mesh limit; the Euclidean Hilbert space of
Section~\ref{ex-reflection} needs the subsequent slab and phase estimates
to support the phase decomposition and the relativistic reconstruction proved
in the later sections; and separated-correlation
compactness does not supply the operator-product remainder or the trace
asymptotic. Until those inputs have detailed proofs or checked applicable
references, this article does not contain a complete existence proof.
